\documentclass[11pt,a4paper]{article}
\usepackage{gtu-el-yazisi}
\begin{document}
\notbilgi{Makine Öğrenmesi}{Metrikler, mesafe ölçüleri, sınıflandırma ölçütleri}{Ders-2.pdf}

\sayfa{1}
\section{Independent component analysis}
\elcizimi[0.45]{sekiller/Ders-2-s1-1.png}[1]

zero shot \quad mesela sınıfın yarısından örnek alındı, diğer tarafta test edildi. diğer taraftan örnek alınmamış.

\section{Metrikler}
\[ \vec{x} = \begin{bmatrix} x_1 \\ x_2 \\ \vdots \\ x_N \end{bmatrix} \]

\subsection{Mesafe ölçümü ve metrikler:}
\textcolor{kirmizikalem}{Distance measure and metrics}

\subsection{Metrik:}
\begin{enumerate}[label=\textcircled{\small\arabic*}]
\item Negatif olmama: $d(x) \ge 0$\footnote{Claude notu: muhtemelen $d(x,y)\ge 0$ olmalı.}
\item $d(x,y)=0 \Rightarrow x=y$
\item $d(x,y)=d(y,x)$ \quad simetri özelliği
\item $d(x,z) \le d(x,y)+d(y,z)$
\end{enumerate}

$\arg\max_x f(x)$, $f(x)$ max yapan

\sayfa{2}
\[ \vec{x} = \begin{bmatrix} x_1 \\ x_2 \\ \vdots \\ x_N \end{bmatrix} \qquad d(x,y) = \sqrt{(x_1-y_1)^2+(x_2-y_2)^2+\cdots} \]

\textbf{Divergence:} (Bu metrik değil ama o da mesafe ölçer) birisi diğerine daha uzak mı onu ölçüyor.

Tüm metrik özelliklerini sağlamazlar

norm:
\[
\left.
\begin{aligned}
\|\vec{x}\|_2 &= \Bigl(\sum_{i=1}^{N} |x_i|^2\Bigr)^{1/2} \\
\|\vec{x}\|_p &= \Bigl(\sum_{i=1}^{N} |x_i|^p\Bigr)^{1/p}
\end{aligned}
\right\} \ell_p\text{-norm}
\]

\elcizimi[0.5]{sekiller/Ders-2-s2-1.png}[2]

\[ \|\vec{x}\|_1 = \sum_{i=1}^{N} |x_i| \]
$\ell_1$-norm: Manhattan norm \\
$\ell_2$-norm: Öklid \\
$\ell_\infty$-norm: $\|x\|_\infty = \max_i |x_i|$

\elcizimi[0.55]{sekiller/Ders-2-s2-2.png}[2]
$\ell_p$-norma göre birim çemberler

\sayfa{3}
\elcizimi[0.5]{sekiller/Ders-2-s3-1.png}[3]
$\rightarrow$ Loss functionlarda normlar kullanılıyor

\[ \|\alpha x\|_p = |\alpha|\,\|x\|_p \]

Okuma alışkanlığını şehirlere göre

\elcizimi[0.45]{sekiller/Ders-2-s3-2.png}[3]
\[ \vec{x}_{izmir} = \begin{bmatrix} x_1 \\ x_2 \\ \vdots \\ x_{14} \end{bmatrix} \]

\elcizimi[0.45]{sekiller/Ders-2-s3-3.png}[3]
burada sıkıntı ne \\
mesela nüfusa bölebiliriz

\[ \|\vec{x}_{izmir}-\vec{x}_{ankara}\|_2 > \|\vec{x}_{izmir}-\vec{x}_{antep}\|_2 \]

\elcizimi[0.4]{sekiller/Ders-2-s3-4.png}[3]
\[ \vec{x}_{izmir}\cdot\vec{x}_{ankara} = \frac{\|\vec{x}_{izmir}\|\,\|\vec{x}_{ankara}\|\cos\alpha}{\|\ \|\cdot\|\ \|} \]

\sayfa{4}
Ama ben olasılık dağılımı üzerinden gitmek istiyorum

\subsection{KL-Divergence:}
\[ D_{KL}(P\|Q) = \sum_{x\in\mathcal{X}} P(x)\cdot\log\frac{P(x)}{Q(x)} \]
bir kez izmir ankara \\
bir kez izmir antep

Ama bu metrik değil

çünkü
\begin{enumerate}[label=\arabic*)]
\item negatif olabilir
\item simetrik değil \quad $d(x,y)\neq d(y,x)$ her zaman
\end{enumerate}

Metriğe çevirelim
\[ D_{JS}(P,Q) = \left\{ \begin{aligned} & \sum_{x\in\mathcal{X}} P(x)\left|\log\frac{P(x)}{Q(x)}\right| \\ & + \sum_{x\in\mathcal{X}} P(x)\left|\log\frac{Q(x)}{P(x)}\right| \end{aligned} \right. \]

\subsection{Jenson-Shannon div}
\[ M(x) = \frac{P(x)+Q(x)}{2} \]
\[ \tfrac12 D_{KL}(P\|M) + \tfrac12 D_{KL}(Q\|M) \]

\elcizimi[0.3]{sekiller/Ders-2-s4-1.png}[4]
Entropi:
\[ H(x) = -\sum_{x\in\mathcal{X}} P(x)\log P(x) \]
\[ P(0)\cdot\log P(0) + P(1)\log P(1) \]
\[ H_1 = -\left(\tfrac23\log\tfrac23 + \tfrac13\log\tfrac13\right) \]
\[ H_2(x) = 0+0 = 0 \]

Ne kadar çok kaos varsa entropi o kadar yükselir.

\sayfa{5}
$P(0)=p$ \qquad $P(1)=1-p$ \qquad $p=0{,}5$ max olur $H(x)$

\subsection{Continuos için:}
\[ H(x) = -\int_{-\infty}^{+\infty} p(x)\log p(x)\,dx \]

\section{Cross-Entropy:}
\[ H(P,Q) = -\sum_{x\in\mathcal{X}} P(x)\log Q(x) \]
$\uparrow$ P ile Q'nun birbirine uzaklığı

\elcizimi[0.75]{sekiller/Ders-2-s5-1.png}[5]

\elcizimi[0.45]{sekiller/Ders-2-s5-2.png}[5]

\[ d_M(x,y) = \sqrt{(x-y)^T\,\underbrace{\Sigma_{xy}^{-1}}_{\text{kovaryans matrisinin tersi}}\,(x-y)} \]

\elcizimi[0.3]{sekiller/Ders-2-s5-3.png}[5]
\[ \Sigma = \begin{bmatrix} 10^{12} & 0 \\ 0 & 10^{4} \end{bmatrix} \]

\sayfa{6}
Bir haber başkasına benziyor mu

Iğdır'da trafik kazasından 10 kişi öldü

$\Big\{$ Tokenization

{\color{turuncukalem}\fbox{\color{black}ığdır trafik kaza 10 kişi ölmek}}

{\color{turuncukalem}\fbox{\color{black}gs fb 3 1 yemek}}

ığdır avokada getirmek

\subsection{Küme Metrikleri}
Hamming Distance \quad Bütün veriler tek (set())
\[ d_H(A,B) = |A \cap B| \]
Jaccard distance:
\[ 1 - \frac{|A\cap B|}{|A\cup B|} \]

\section{Sınıflandırma metrikleri}
\elcizimi[0.4]{sekiller/Ders-2-s6-1.png}[6]

\subsection{Doğruluk (Acc)}
\[ \frac{TN+TP}{\#\text{total}} = Acc \]
\[ \{x^t, r^t\}_{t=1}^{N} \]
\[ r^t = 0,\ t=1{:}N \]
\[ \text{Precision} = \frac{TP}{TP+FP} \]

\sayfa{7}
\[ \text{Recell} = \frac{TP}{TP+FN} \]
\[ \text{F-1 score} = 2\times\frac{\text{Prec}\times\text{Recell}}{\text{Prec}+\text{Recell}} \]

\elcizimi[0.55]{sekiller/Ders-2-s7-1.png}[7]

\elcizimi[0.55]{sekiller/Ders-2-s7-2.png}[7]

\elcizimi[0.5]{sekiller/Ders-2-s7-3.png}[7]
\[ FPR = \frac{FP}{TN+FP} \]

0'ı eşik değer yaparsam

\elcizimi[0.3]{sekiller/Ders-2-s7-4.png}[7]
\[ TPR = \frac{4}{0+4} = 1 \qquad FPR = \frac{5}{5+0} = 1 \]

\elcizimi[0.55]{sekiller/Ders-2-s7-5.png}[7]

\elcizimi[0.3]{sekiller/Ders-2-s7-6.png}[7]
\[ TPR = \frac{4}{4+0} = 1 \qquad FPR = \frac{3}{3+2} = \frac{3}{5} \]

\sayfa{8}
Soru

\elcizimi[0.6]{sekiller/Ders-2-s8-1.png}[8]

\elcizimi[0.45]{sekiller/Ders-2-s8-2.png}[8]

\elcizimi[0.3]{sekiller/Ders-2-s8-3.png}[8]
\[ TPR = \frac{2}{2+1} = \frac{2}{3} \qquad FPR = \frac{2}{2+3} = \frac{2}{5} \]

\sayfa{9}
\section{Confusion matrix}
\elcizimi[0.5]{sekiller/Ders-2-s9-1.png}[9]
\[ \{x^t, r^t\}_{t=1}^{N} \]
\[ r^t \in C \]
\[ r^t \in \mathbb{R}^N \]

\elcizimi[0.3]{sekiller/Ders-2-s9-2.png}[9]
\[ f(x^t) = r^{\alpha} \]
\[ \frac{1}{N}\sum_{t=1}^{N}\bigl(\underbrace{\hat r^t}_{\text{tahminler}} - r^t\bigr)^2 \]

\elcizimi[0.55]{sekiller/Ders-2-s9-3.png}[9]
\[ ax+b = f(x) \]

\end{document}
